\section*{Introduction}

Consider scattering $n$ points around a circle by repeatedly stepping by a fixed
angle $\alpha$. The points land at fractional parts
\[
  \{1\cdot\alpha\},\; \{2\cdot\alpha\},\; \ldots,\; \{n\cdot\alpha\}
\]
where $\{x\} = x - \lfloor x \rfloor$ denotes the fractional part of $x$, and we
identify $[0,1)$ with the unit circle.

If $\alpha$ is rational, say $\alpha = p/q$ in lowest terms, the sequence is
eventually periodic with period $q$ and the points cluster into $q$ equally-spaced
positions. The more interesting case is \emph{irrational} $\alpha$, where no two
points ever coincide and the orbit is everywhere dense in $[0,1)$ by Weyl's
equidistribution theorem.

\medskip
\noindent\textbf{The question.} How many distinct gap lengths appear between
consecutive points?

\medskip
Hugo Steinhaus posed this as a conjecture around 1958: no matter which irrational
$\alpha$ you choose and no matter how many points $n$ you place, the $n$ gaps
between consecutive points take \emph{at most three} distinct lengths. Moreover,
when there are three lengths, the largest equals the sum of the other two.

The result was proved independently and almost simultaneously by
Vera~T.\ S\'os (1958), Jan Suranyi, and Stanisław Świerczkowski. It is now
known as the \textbf{Three-Distance Theorem} (or Steinhaus theorem, or
three-gap theorem).
